% =============================================================================
\section{Extenders}
\label{sec:extenders}
% =============================================================================
In several places in our binding code we are given a type that models
certain concepts and based on several conditions, this type must be replaced 
with a new type that models additional concepts. Instead of using old style 
code in the form of \cppLstinline{#if,#elif,#else,#endif} directives, we 
exploit generic programming and implement these extenders using modern code 
that is type safe. In the following we present four such scenarios.

% -----------------------------------------------------------------------------
\subsection{2D Arrangement Traits Extender}
\label{sec:extenders:arrTraits}
% -----------------------------------------------------------------------------
When the user enables the generation of bindings for the 
\iiDArrangementsPackage{} package, the binding for an instance of the class
template \cppLstinline{Arrangement_2<Traits, Dcel>} is generated. When this
class template is instantiated the \cppLstinline{Traits} template parameter 
must be substituted with a specific traits model. By default an instance of 
the \cppLstinline{Arr_segment_traits_2<Kernel>} is used; however, the user 
can override this selection and choose any one of the supported geometry 
traits. We refer to this traits as the basic geometry traits. If the user 
also enables the generation of bindings for the 
\iiDRegularizedBooleanSetOperationsPackage{} package, the basic traits must
be extended, because it must also model the 
\cppLstinline{GeneralPolygonSetTraits_2} concept 
(see Figure~\ref{fig:atc:gps}); see Section~\ref{ssec:modules:aos2}. The
extension however differs for different basic traits. It is automatically
done via a dedicated class template and four specializations shown in Listing~\ref{lst:arrTraitsExtender}. The final traits is obtained
using this extender as shown in Listing~\ref{lst:arrFinalTraits}. Here,
\cppLstinline{boolean_set_operations_2_bindings()} is a constant binary
function (evaluated at compile time) that determines whether the generation 
of bindings for the 
\iiDRegularizedBooleanSetOperationsPackage{} package is enabled, and
\cppLstinline{BGT} is the basic geometry traits.

\begin{lstfloat}
  \begin{lstlisting}[style=cppFramedStyle,basicstyle=\normalsize,%
                     label={lst:arrTraitsExtender},%
                     caption={2D Arrangement Traits Extender.}]
// Traits extender
template <bool b, typename Base> struct Tr {};

// Specialization for the case where bindings for Bso2 are not generated
template <typename Base>
struct Tr<false, Base> { typedef Base type; };

// Specialization for the general case where bindings for Bso2 are generated
template <typename Base>
struct Tr<true, Base> { typedef CGAL::Gps_traits_2<Base> type; };

// Specialization for the segment traits
template <>
struct Tr<true, CGAL::Arr_segment_traits_2<Kernel>> {
  typedef CGAL::Gps_segment_traits_2<Kernel, Point_2_container> type;
};

// Specialization for the circle segment traits
template <>
struct Tr<true, CGAL::Arr_circle_segment_traits_2<Kernel>> {
  typedef CGAL::Gps_circle_segment_traits_2<Kernel> type;
};
  \end{lstlisting}
\end{lstfloat}

\begin{lstfloat}
  \begin{lstlisting}[style=cppFramedStyle,basicstyle=\normalsize,%
                     label={lst:arrFinalTraits},%
                     caption={The final 2D Arrangement traits type definition.}]
typedef Tr<boolean_set_operations_2_bindings(), BGT>::type      Geometry_traits;
  \end{lstlisting}
\end{lstfloat}

% -----------------------------------------------------------------------------
\subsection{2D Arrangement DCEL Cell Extenders}
\label{sec:extenders:arrDcel}
% -----------------------------------------------------------------------------
When the user enables the generation of bindings for the 
\iiDArrangementsPackage{} package, the binding for an instance of the class
template \cppLstinline{Arrangement_2<Traits, Dcel>} is generated. When this
class template is instantiated the \cppLstinline{Dcel} template-parameter 
should be substituted with a type that models the 
\cppLstinline{ArrangementDcel} concept. We substitute this parameter with an
instance of the template \cppLstinline{CGAL::Arr_dcel_base<V, H, F>}, 
where \cppLstinline{V}, \cppLstinline{H}, and \cppLstinline{F} must  
model the concepts \cppLstinline{ArrangementDcelVertex},
\cppLstinline{ArrangementDcelHalfedge}, and 
\cppLstinline{ArrangementDcelFace}, respectively. By default, we substitute 
these parameters with the instances \cppLstinline{Arr_vertex_base<Traits::Point_2>},
\cppLstinline{Arr_halfedge_base<Traits::X_monotone_curve_2>}, and
\cppLstinline{Arr_face_base}, respectively, referred to as the basic types
of cells. If the user chooses to extend the vertex type, then we substitute 
\cppLstinline{V} with the instance 
\cppLstinline{Arr_extended_ vertex<Vb, bp::object>}, where \cppLstinline{Vb}
is the basic vertex type above. This conditional extension is carried out via 
a dedicated class template and two specializations shown in 
Listing~\ref{lst:arrDcelVertexExtender}. The final vertex type is obtained
using this extender as shown in Listing~\ref{lst:arrFinalVertex}. Here,
\cppLstinline{is_vertex_extended()} is a constant binary
function that determines whether the user extension of vertices is enabled.

\begin{lstfloat}
  \begin{lstlisting}[style=cppFramedStyle,basicstyle=\normalsize,%
                     label={lst:arrDcelVertexExtender},%
                     caption={2D Arrangement DCEL user extended vertex type extender.}]
// Vertex extender
template <bool b, typename Vb> struct Vertex_extender {};

// Specialization
template <typename Vb> struct Vertex_extender<false, Vb> { typedef Vb type; };

// Specialization
template <typename Vb> struct Vertex_extender<true> 
{ typedef CGAL::Arr_extended_vertex<Vb, bp::object> type; };
  \end{lstlisting}
\end{lstfloat}

\begin{lstfloat}
  \begin{lstlisting}[style=cppFramedStyle,basicstyle=\normalsize,%
                     label={lst:arrFinalVertex},%
                     caption={The final 2D Arrangement vertex type definition.}]
CGAL::Arr_vertex_base<Geometry_traits::Point_2>                         Vb;
typedef Vertex_extended<is_vertex_extended(), Vb, bp::object>::type     V;
  \end{lstlisting}
\end{lstfloat}

If the user also enables the generation of bindings for the 
\iiDRegularizedBooleanSetOperationsPackage{} package, the template 
parameters \cppLstinline{H} and \cppLstinline{F} must be substituted with 
types that also model the  
\cppLstinline{GeneralPolygonSetDcelHalfedge} and
\cppLstinline{GeneralPolygonSetDcelFace} concepts, respectively; see 
Section~\ref{ssec:modules:bso2}. In addition, similar to the case of
vertex extension, if the user chooses to extend the halfedge or face types, 
then we substitute \cppLstinline{V} with the instances
\cppLstinline{Arr_extended_halfedge<Hb, bp::object>} or
\cppLstinline{Arr_extended_ face<Fb, bp::object>}, respectively, where
\cppLstinline{Hb} and \cppLstinline{Fb} are the types above. The 
conditional extensions of the halfedge type are carried out via dedicated 
class templates and corresponding specializations shown in 
Listing~\ref{lst:arrDcelBasicHalfedgeExtender} and
Listing~\ref{lst:arrDcelHalfedgeExtender}. The final halfedge type is obtained
using these extenders as shown in Listing~\ref{lst:arrFinalHalfedge}. Here,
\cppLstinline{boolean_set_operations_2_bindings()} is a constant binary
function that determines whether the generation of bindings for the 
\iiDRegularizedBooleanSetOperationsPackage{} package is enabled, and
\cppLstinline{is_halfedge_extended()} is a constant binary function that 
determines whether the user extension of halfedges is enabled.

\begin{lstfloat}
  \begin{lstlisting}[style=cppFramedStyle,basicstyle=\normalsize,%
                     label={lst:arrDcelBasicHalfedgeExtender},%
                     caption={2D Arrangement DCEL Halfedge extender for 2D Regularized Boolean set operations.}]
// Basic halfedge extender
template <bool b> struct Halfedge_gps {};

// Specialization
template <> struct Halfedge_gps<false>
{ typedef CGAL::Arr_halfedge_base<Geometry_traits::X_monotone_curve_2> type; };

// Specialization
template <> struct Halfedge_gps<true>
{ typedef CGAL::Gps_halfedge_base<Geometry_traits::X_monotone_curve_2> type; };
  \end{lstlisting}
\end{lstfloat}

\begin{lstfloat}
  \begin{lstlisting}[style=cppFramedStyle,basicstyle=\normalsize,%
                     label={lst:arrDcelHalfedgeExtender},%
                     caption={2D Arrangement DCEL user extended halfedge type extender.}]
// Halfedge extender
template <bool b, typename Hb> struct Halfedge_extender {};

// Specialization
template <typename Hb> struct Halfedge_extender<false, Hb> { typedef Hb type; };

// Specialization
template <typename Hb> struct Halfedge_extender<true, Hb> 
{ typedef CGAL::Arr_extended_halfedge<Hb, bp::object> type; };
  \end{lstlisting}
\end{lstfloat}

\begin{lstfloat}
  \begin{lstlisting}[style=cppFramedStyle,basicstyle=\normalsize,%
                     label={lst:arrFinalHalfedge},%
                     caption={The final 2D Arrangement DCEL halfedge type definition.}]
typedef Halfedge_gps<boolean_set_operations_2_bindings()>::type         Hb;
typedef Halfedge_extended<is_halfedge_extended(), Hb, bp::object>::type H;
  \end{lstlisting}
\end{lstfloat}

Similarly, the conditional extensions of the face type are carried out 
via dedicated class templates and corresponding specializations shown in 
Listing~\ref{lst:arrDcelBasicFaceExtender} and
Listing~\ref{lst:arrDcelFaceExtender}. The final face type is obtained
using these extenders as shown in Listing~\ref{lst:arrFinalFace}.
\begin{lstfloat}
  \begin{lstlisting}[style=cppFramedStyle,basicstyle=\normalsize,%
                     label={lst:arrDcelBasicFaceExtender},%
                     caption={2D Arrangement DCEL Face extender for 2D Regularized Boolean set operations.}]
// Basic face extender
template <bool b> struct Face_gps {};

// Specialization
template <> struct Face_gps<false> { typedef CGAL::Arr_face_base type; };

// Specialization
template <> struct Face_gps<true> { typedef CGAL::Gps_face_base type; };
  \end{lstlisting}
\end{lstfloat}

\begin{lstfloat}
  \begin{lstlisting}[style=cppFramedStyle,basicstyle=\normalsize,%
                     label={lst:arrDcelFaceExtender},%
                     caption={2D Arrangement DCEL user extended face type extender.}]
// Face extender
template <bool b, typename Fb> struct Face_extender {};

// Specialization
template <typename Fb> struct Face_extender<false, Fb> { typedef Fb type; };

// Specialization
template <typename Fb> struct Face_extender<true, Fb> 
{ typedef CGAL::Arr_extended_face<Fb, bp::object> type; };
  \end{lstlisting}
\end{lstfloat}

\begin{lstfloat}
  \begin{lstlisting}[style=cppFramedStyle,basicstyle=\normalsize,%
                     label={lst:arrFinalFace},%
                     caption={The final 2D Arrangement DCEL face type definition.}]
typedef Face_gps<boolean_set_operations_2_bindings()>::type             Fb;
typedef Face_extended<is_face_extended(), Fb, bp::object>::type         F;
  \end{lstlisting}
\end{lstfloat}

% -----------------------------------------------------------------------------
\subsection{2D Triangulation Cell Extenders}
\label{sec:extenders:iitri}
% -----------------------------------------------------------------------------
When the user enables the generation of bindings for the 
\iiDTriangulationsPackage{} package, the binding for an instance of one
of the 2D triangulation class templates is generated. Each of these
class templates has a template parameter that must be substituted with 
a model of the concept \cppLstinline{TriangulationDataStructure_2} when
the triangulation class template is instantiated. We substitute this
parameter with an instance of the class template 
\cppLstinline{Triangulation_data_structure_2<V, F>}. 

The type that substitutes the \cppLstinline{V} template parameter when 
the class template \cppLstinline{Triangulation_data_structure_ 2<V, F>} 
is instantiated must model the \cppLstinline{TriangulationDSVertexBase_2}
concepts. If the binding is generated for a periodic triangulation, this
parameter must be substituted with a type that also models the concept
\cppLstinline{Periodic_2TriangulationDSVertexBase_2}. The extension is 
automatically done via a dedicated class template and three 
specializations shown in Listing~\ref{lst:iitriBaseVertexExtender}. Here 
\cppLstinline{Traits} is the geometry traits type; see 
Section~\ref{ssec:modules:tri2} for an explanation on how this type is 
determined.

\begin{lstfloat}
  \begin{lstlisting}[style=cppFramedStyle,basicstyle=\normalsize,%
                     label={lst:iitriBaseVertexExtender},%
                     caption={The 2D Triangulation base vertex type extender.}]
// Vertex base selection
template <bool b, int i> struct Vertex_base_name {};

// Specialization
template <int i> struct Vertex_base_name<false, i>
{ typedef CGAL::Triangulation_vertex_base_2<Traits> type; };

// Specialization
template <> struct Vertex_base_name<false, CGALPY_TRI2_REGULAR>
{ typedef CGAL::Regular_triangulation_vertex_base_2<Traits> type; };

// Specialization
template <int i> struct Vertex_base_name<true, i>
{ typedef CGAL::Periodic_2_triangulation_vertex_base_2<Traits> type; };
  \end{lstlisting}
\end{lstfloat}

If the user chooses to extend the vertex type, the type that substitutes
the \cppLstinline{V} template parameter is further extended. This is
automatically done via yet another dedicated class template and two
specializations shown in Listing~\ref{lst:iitriExtendedVertexExtender}. 

\begin{lstfloat}
  \begin{lstlisting}[style=cppFramedStyle,basicstyle=\normalsize,%
                     label={lst:iitriExtendedVertexExtender},%
                     caption={The 2D Triangulation user extended vertex type extender.}]
// Vertex with info
template <bool b, typename Vb> struct Vertex_with_info {};

// Specialization
template <typename Vb> struct Vertex_with_info<false, Vb> { typedef Vb type; };

// Specialization
template <typename Vb> struct Vertex_with_info<true, Vb>
{ typedef CGAL::Triangulation_vertex_base_with_info_2<bp::object, Traits, Vb> type; };
  \end{lstlisting}
\end{lstfloat}

If the triangulation is used to define an alpha shape type, the template
parameter \cppLstinline{V} must be substituted with a type that also models
the concept \cppLstinline{AlphaShapeVertex_2}; see 
Section~\ref{ssec:modules:as2}. This is automatically done via yet another 
class template shown in Listing~\ref{lst:iitriASVertexExtender}.

\begin{lstfloat}
  \begin{lstlisting}[style=cppFramedStyle,basicstyle=\normalsize,%
                     label={lst:iitriASVertexExtender},%
                     caption={The 2D Triangulation for alpha shapes vertex type extender.}]
// Vertex alpha shape
template <bool b, typename Vb, typename ExactComparison> 
struct Vertex_alpha_shape {};

// Specialization
template <typename Vb, typename ExactComparison>
struct Vertex_alpha_shape<false, Vb, ExactComparison> { typedef Vb type; };

// Specialization
template <typename Vb, typename ExactComparison>
struct Vertex_alpha_shape<true, Vb, ExactComparison>
{ typedef CGAL::Alpha_shape_vertex_base_2<Traits, Vb, ExactComparison> type; };
  \end{lstlisting}
\end{lstfloat}

Finally, if the binding is generated for a hierarchy triangulation, e.g., 
an instance of the template parameter
\cppLstinline{Triangulation_hierarchy_2<Triangulation_2>}, the
\cppLstinline{V} template parameter must be substituted with
a type that also models the concept
\cppLstinline{TriangulationHierarchyVertexBase_2}. (Observe that
there is no special requirements on the type that substitutes the
\cppLstinline{F} parameter in this case.) This is automatically done via
yet another class template shown in 
Listing~\ref{lst:iitriHierarchyVertexExtender}.

\begin{lstfloat}
  \begin{lstlisting}[style=cppFramedStyle,basicstyle=\normalsize,%
                     label={lst:iitriHierarchyVertexExtender},%
                     caption={The 2D Triangulation hierarchy vertex type extender.}]
// Vertex triangulation hierarchy
template <bool b, typename Vb> struct Vertex_hierarchy {};

// Specialization
template <typename Vb> struct Vertex_hierarchy<false, Vb> { typedef Vb type; };

// Specialization
template <typename Vb> struct Vertex_hierarchy<true, Vb>
{ typedef CGAL::Triangulation_hierarchy_vertex_base_2<Vb> type; };
  \end{lstlisting}
\end{lstfloat}

% Basic
%%%%%%%
The type \cppLstinline{Vb} is defined as the vertex base type.  If
bindings are generated for a periodic triangulation, the type
\cppLstinline{Vb} is defined as 
\cppLstinline{Periodic_2_triangulation_vertex_base_2<Traits>};
otherwise, if bindings are generated for an instance of
\cppLstinline{Regular_triangulation_2<Traits, Tds>}, the type 
\cppLstinline{Vb} is defined as
\cppLstinline{Regular_triangulation_ vertex_base_2<Traits>}; otherwise, 
the type \cppLstinline{Vb} is defined as
\cppLstinline{Triangulation_vertex_ base_2<Traits>}; 
% Extended
%%%%%%%%%%
If the user enables vertex-type extension, the type \cppLstinline{Vbi} 
is defined as
\cppLstinline{Triangulation_vertex_base_with_ info_2<bp::object, Kernel, Vb>}; 
otherwise the type \cppLstinline{Vbi} is defined as \cppLstinline{Vb}.
% Alpha Shape
%%%%%%%%%%%%%
If the triangulation is used to define an alpha shape type, the type 
\cppLstinline{Vbia} is defined as
\cppLstinline{Alpha_shape_vertex_base_2<Traits, Vb, ExactComparison>}, 
where \cppLstinline{ExactComparison} is substituted with either 
\cppLstinline{true} or \cppLstinline{false} based on the setting of the 
\cmake{} flag \cmakeLstinline{CGALPY_AS2_EXACT_COMPARISON}; see
Table~\ref{tab:as2}. Otherwise, \cppLstinline{V} is defined as
\cppLstinline{Vbi}.
% Hierarchy
%%%%%%%%%%%
If the binding is generated for a hierarchy triangulation, the final 
vertex type \cppLstinline{V} is defined as
\cppLstinline{Triangulation_hierarchy_vertex_ base_2<Vbi>}; otherwise
the final vertex type \cppLstinline{V} is defined as \cppLstinline{Vbia};
see Listing~\ref{lst:iitriFinalVertex}. Here, 
\cppLstinline{is_periodic()},
\cppLstinline{vertex_with_ info()},
\cppLstinline{alpha_shape_2_bingings()}, and
\cppLstinline{hierarchy()} are constant binary functions.
\cppLstinline{is_periodic()} that determines whether the triangulation, 
binding for which is generated, is periodic.
\cppLstinline{vertex_with_info()} determines whether the user 
extension of vertices is enabled.
\cppLstinline{alpha_shape_2_bingings()} indicates whether bindings for
the \iiDAlphaShapesPackage{} package.
\cppLstinline{hierarchy()} indicates whether the binding is generated
for a hierarchy triangulation. 

\begin{lstfloat}
  \begin{lstlisting}[style=cppFramedStyle,basicstyle=\normalsize,%
                     label={lst:iitriFinalVertex},%
                     caption={The final 2D Triangulation vertex type definition.}]
typedef Vertex_base_name<is_periodic(), CGALPY_TRI2>::type           Vb;
typedef Vertex_with_info<vertex_with_info(), Vb>::type               Vbi;
typedef Vertex_alpha_shape<alpha_shape_2_bindings(), Vbi, Ec>::type  Vbia;
typedef Vertex_hierarchy<hierarchy(), Vbia>::type                    V;
  \end{lstlisting}
\end{lstfloat}

The final face type is determined in a similar fashion. The only
difference between the selections of the vertex and face types is that
the setting of the \cmake{} flag \cmakeLstinline{CGALPY_TRI2_HIERARCHY} does
not have an effect on the face type selection.

% -----------------------------------------------------------------------------
\subsection{3D Triangulation Cell Extenders}
\label{sec:extenders:iiitri}
% -----------------------------------------------------------------------------
When the user enables the generation of bindings for the 
\iiiDTriangulationsPackage{} package, the binding for an instance 
of one of the 3D triangulation class templates is generated; see 
Section~\ref{ssec:modules:tri3}. The scenario here is similar to
the scenario of the \iiDTriangulationsPackage{} package described
in Section~\ref{sec:extenders:iitri}. We skip the description of
the extenders and jump to the description of their application.

If bindings are generated for a periodic triangulation, the type
\cppLstinline{Vbp} is defined as
\cppLstinline{Periodic_3_triangulation_ ds_vertex_base_3<>};
otherwise, \cppLstinline{Vbp} is defined as
\cppLstinline{Triangulation_ds_vertex_base_3<>}.
%
If bindings are generated for an instance of either
\cppLstinline{Regular_triangulation_vertex_base_3<Kernel>} or
\cppLstinline{Periodic_3_ Regular_triangulation_vertex_base_3<Kernel>},
the type \cppLstinline{Vb} is defined as
\cppLstinline{Regular_triangulation_ vertex_base_3<Traits, Vbp>};
otherwise, the type \cppLstinline{Vb} is defined as
\cppLstinline{Triangulation_vertex_base_3<Traits, Vbp>}.
%
If the vertex type is extended, the type \cppLstinline{Vbi} is defined 
as 
\cppLstinline{Triangulation_vertex_base_ with_info_3<bp::object, Traits, Vb>};
otherwise, the type \cppLstinline{Vbi} is defined as \cppLstinline{Vb}.
%
If the binding is generated for a hierarchy triangulation, the type 
\cppLstinline{Vbih} is defined as
\cppLstinline{Triangulation_hierarchy_vertex_base_3<Vbi>}; otherwise
the type \cppLstinline{Vbih} is defined as \cppLstinline{Vbi}.
%
If the triangulation is used to define an alpha shape type,
the final type \cppLstinline{V} (that substitutes the
\cppLstinline{V} parameter) is defined as either
\cppLstinline{Alpha_shape_vertex_base_3<Traits, Vbih, ExactComparison>} 
or
\cppLstinline{Fixed_Alpha_shape_vertex_base_3<Traits, Vbih, ExactComparison>} 
based on the selection of the alpha shape class (either fixed or plain); 
see Table~\ref{tab:as3}; \cppLstinline{ExactComparison} is substituted with 
either \cppLstinline{true} or \cppLstinline{false} based on the user 
selection of exact comparisons; see Table~\ref{tab:as3}. Otherwise, 
\cppLstinline{V} is defined as \cppLstinline{Vbih}; see 
Listing~\ref{lst:iiitriFinalVertex}. (Observe that while 3D periodic 
regular triangulations are supported in the form of the class template 
\cppLstinline{Periodic_3_regular_triangulation_3},
corresponding 2D triangulations are not; thus the difference between
the structures of the code in Listing~\ref{lst:iitriFinalVertex} and Listing~\ref{lst:iiitriFinalVertex}.

\begin{lstfloat}
  \begin{lstlisting}[style=cppFramedStyle,basicstyle=\normalsize,%
                     label={lst:iiitriFinalVertex},%
                     caption={The final 3D Triangulation vertex type definition.}]
typedef Vertex_periodic<is_periodic()>::type                         Vb;
typedef Vertex_regular<is_regular(), Vb>::type                       Vbr;
typedef Vertex_with_info<vertex_with_info(), Vbr>::type              Vbri;
typedef Vertex_alpha_shape<alpha_shape_3_bindings(), Vbri, Ec>::type Vbria;
typedef Vertex_hierarchy<hierarchy(), Vbria>::type                   V;
  \end{lstlisting}
\end{lstfloat}

The final cell type is determined in a similar fashion. The only
difference between the selections of the vertex and cell types is that
the setting of the \cmake{} flag \cmakeLstinline{CGALPY_TRI3_HIERARCHY}
does not have an effect on the cell type selection. 
